\documentclass{situationist-preprint}

\title{Why Promises Bind}
\msid{ETH-16-a}
\msfamily{ETH-16}
\msdiscipline{Normative and Applied Ethics}
\mstype{Principal manuscript}
\mskeywords{promising; promissory obligation; reliance; assurance; Scanlon; conventionalism; normative powers; release}

\begin{document}
\maketitle

\begin{abstract}
A promise creates an obligation where there was none. Conventionalists locate the source of the obligation in a social practice; will theorists, in the promisor's exercise of a normative power; expectation theorists, in the expectations the promise creates. The expectation view is the most explanatory, but it faces a dilemma: if the obligation rests on the promisee's actual reliance, it lapses when the promisee forgets, and if it does not, it is unclear what reliance has to do with it. We argue that promissory obligation arises from the creation of reasonable reliance, and that it persists even when the promisee forgets. The key move is to distinguish the promisee's entitlement to rely, which the promisor's act creates, from her psychological state of relying, which may come and go. The obligation is grounded in a past act and directed to a person's entitlement, and such obligations persist until they are discharged or released. Forgetting is neither. The account explains why insincere promises bind, why uncommunicated resolutions do not, and why only the promisee can release the promisor.
\end{abstract}

\section{The problem}

Hume held that the obligation of promises is artificial: a promise is a form of words whose obligation is invented by convention, because it is in everyone's interest that there be a way of binding oneself to future performance \citep[3.2.5]{Hume1739}. \citet{Rawls1955} gave the conventionalist view its modern form, on which promising is a practice defined by constitutive rules, and the obligation to keep a particular promise derives from the duty of fairness not to take advantage of a practice from which one benefits. Against conventionalism, \citet{Scanlon1990,Scanlon1998} argued that the wrong of breaking a promise is a special case of a wrong that requires no convention: the wrong of leading someone to form expectations about what one will do, when she wants assurance, and then failing to act accordingly. His Principle of Fidelity says, roughly, that if one has voluntarily and intentionally led another to expect that one will do something, knowing that she wants to be assured of it, one must do it unless she consents otherwise. A third family of views holds that a promise is the exercise of a normative power to create an obligation by communicating the intention to create it \citep{Raz1977,Owens2012}.

The expectation view faces two problem cases. In the first, the promisee never relies on the promise: she does not believe it, or does not change her plans on account of it. In the second, she forgets it. Suppose I promise a colleague that I will read her draft by Friday, and on Thursday she has forgotten I ever offered. If the obligation rested on her reliance, it would now have lapsed, and I could skip the reading with a clear conscience. Hardly anyone believes this. \citet{KolodnyWallace2003} argued that cases of this kind push the expectation theorist back towards a practice, since only a practice can explain an obligation that outlasts the expectations it was supposed to protect.

We argue that the expectation view, correctly stated, handles both cases. Promissory obligation arises from the creation of reasonable reliance, and it persists even when the promisee forgets. The first problem case is treated in the companion manuscript ETH-16-b; the second is our main concern here.

\section{Framework}

\begin{definition}[Promise]\label{d:promise}
A \emph{promise} by $A$ to $B$ to do $X$ is a communicative act by which $A$ intends to lead $B$ to expect that $A$ will do $X$, intends $B$ to recognize this intention, and presents himself as undertaking to do $X$, and which $B$ understands as such.
\end{definition}

What makes an utterance a promise, as a speech act, is a matter of the constitutive rules that individuate it \citep{Searle1969}, as LNG-13 argues for speech acts generally. The question here is different: not what makes an act a promise, but what makes a promise bind.

\begin{definition}[Reasonable reliance]\label{d:reasonable}
$B$ has \emph{reasonable reliance} on $A$'s doing $X$ when $B$ is entitled, in virtue of what $A$ has done, to expect that $A$ will do $X$ and to plan on that expectation. \emph{Actual reliance} is the psychological state of expecting that $A$ will do $X$ and planning on it.
\end{definition}

Reasonable reliance is a normative status, not a psychological state. It can exist without actual reliance, as when the promisee is entitled to plan on the promise but chooses not to, and actual reliance can exist without it, as when someone overhears a resolution not addressed to him and plans on it.

\begin{principle}[Reliance]\label{p:reliance}
If $A$ voluntarily and intentionally creates in $B$ reasonable reliance on $A$'s doing $X$, knowing that $B$ wants assurance that $A$ will do $X$, then $A$ is obligated to $B$ to do $X$, unless $B$ releases him or the obligation is defeated by weightier reasons.
\end{principle}

\begin{principle}[Historical grounding]\label{p:history}
An obligation grounded in a past act of the obligated agent, and directed to another person's entitlement, persists until it is discharged, released by the person to whom it is owed, or defeated; it does not lapse when the psychological states that the act produced in that person cease.
\end{principle}

\begin{assumption}[Release]\label{a:release}
Release is an act of the promisee's will, communicated to the promisor. Forgetting is not release.
\end{assumption}

Principle~\ref{p:reliance} is a version of Scanlon's Principle of Fidelity, restated in terms of reasonable rather than actual reliance; the restatement is what the second problem case requires. Principle~\ref{p:history} states a feature common to the obligations that arise from transactions. A debt persists when the creditor forgets the loan; a duty to compensate persists when the victim has stopped thinking about the injury. Assumption~\ref{a:release} records the platitude that a promisee who forgets has not waived anything: waiver is something one does, and one cannot do it unawares.

\section{The argument}

\begin{lemma}\label{l:creates}
Every promise creates reasonable reliance in the promisee.
\end{lemma}

\begin{proof}
By Definition~\ref{d:promise}, the promisor intends to lead the promisee to expect his performance, intends her to recognize that intention, and presents himself as undertaking to perform. By these acts he invites her to rely on him, and he makes himself answerable to her if she does. An invitation to rely, made by someone who presents himself as answerable for its being well founded, entitles the person invited to rely. By Definition~\ref{d:reasonable}, the promisee has reasonable reliance.
\end{proof}

\begin{theorem}[Source]\label{t:source}
Promissory obligation arises from the creation of reasonable reliance.
\end{theorem}

\begin{proof}
A promisor acts voluntarily and intentionally (Definition~\ref{d:promise}), and in the central case the promisee wants assurance, which is why promises are asked for and given. By Lemma~\ref{l:creates}, the promise creates reasonable reliance. By Principle~\ref{p:reliance}, the promisor is obligated to the promisee to perform. The obligation is derived from the creation of reasonable reliance and from nothing else: no premise about a practice, and none about the promisee's actual reliance, has been used.
\end{proof}

\begin{theorem}[Persistence]\label{t:persist}
Promissory obligation persists when the promisee forgets the promise.
\end{theorem}

\begin{proof}
By Theorem~\ref{t:source}, the obligation is grounded in a past act of the promisor, the creation of reasonable reliance, and it is directed to the promisee's entitlement to rely. Forgetting removes the promisee's actual reliance, a psychological state that the act produced, but it neither discharges the obligation, since the promisor has not performed, nor releases it, by Assumption~\ref{a:release}. By Principle~\ref{p:history}, the obligation persists.
\end{proof}

The proof turns on the difference between the entitlement and the state. The promisee who forgets is still entitled to rely: were she reminded, she could pick up her plans where she left them, and she would have a complaint against a promisor who answered that, since she had forgotten, he had felt free. The complaint is the same complaint she would have had if she had remembered, and it is addressed to the same thing, the falsification of an assurance that she was entitled to have. Forgetting changes what she is thinking about, not what she is owed.

The theorem also explains three further features of promising. \emph{Insincere promises bind}: the promisor who never intends to perform still creates reasonable reliance, so Principle~\ref{p:reliance} applies to him, whereas a view on which the obligation is created by the intention to bind oneself must explain why the absence of that intention does not void the promise. \emph{Uncommunicated resolutions do not bind}: a resolution that I make in private, however firm, creates no entitlement in anyone, and so no promissory obligation, which is why Definition~\ref{d:promise} requires uptake. \emph{Only the promisee can release}: since the obligation is directed to her entitlement, she alone can waive it, and a third party who learns of the promise cannot.

\section{Objections}

\begin{objection}[Circularity]
The promisee is entitled to rely on the promise only because the promisor is obligated to perform. To ground the obligation in the entitlement is to argue in a circle \citep{KolodnyWallace2003}.
\end{objection}

\begin{reply}
The circle would be vicious if the entitlement were derived from the obligation and the obligation from the entitlement. In the proof of Lemma~\ref{l:creates}, the entitlement is derived from the promisor's act of inviting reliance while presenting himself as answerable for it, and not from the obligation. What the promisor is answerable for is the wrong Scanlon identifies, that of leading someone to form expectations she wants to form and then defeating them; that wrong is independent of promising, since it can be done by assurances that are not promises \citep{Scanlon1998}. The obligation and the entitlement are two aspects of a single normative relation created by the act, and neither is prior to the other in the order of grounds; both are grounded in the act.
\end{reply}

\begin{objection}[No interest at stake]
The point of protecting reliance is to protect the promisee's interest in being able to plan. A promisee who has forgotten has no such interest at stake. Breaking the promise sets back no interest of hers, and there is nothing for the obligation to protect.
\end{objection}

\begin{reply}
The interest is not only in the plans she happens to have made, but in being able to plan on the promise at all, and in having her entitlement to do so respected. Both survive forgetting. A promisee may recall the promise at the last moment and need it; and even if she never recalls it, the promisor who breaks it has treated her entitlement as conditional on her attention, which is not how she was entitled to have it treated. \citet{Shiffrin2008} argues that promising serves the interest in relationships of mutual trust between equals; a promise that binds only while the promisee is watching is not a form of trust.
\end{reply}

\begin{objection}[Convention]
Reliance on a mere statement of intention is not reasonable. Only the existence of a practice of promising, with its constitutive rules, makes it reasonable to rely on someone's word, and so the practice, not the reliance, is the source of the obligation \citep{Rawls1955}.
\end{objection}

\begin{reply}
A practice may supply the means of creating reasonable reliance, a recognized form of words; it is not the source of the entitlement, since reasonable reliance can be created without it. Scanlon argues that the wrong of defeating expectations one has deliberately created can be done in circumstances that no social practice governs, and that it is condemned by principles no one could reasonably reject without appeal to any convention \citep{Scanlon1990}. What matters is that the agent invites reliance and presents himself as answerable for it. Where a practice exists, it makes the invitation easier to issue and to recognize, which is why promising is so efficient a way of creating obligations. But the obligation is created by the invitation, as Theorem~\ref{t:source} says.
\end{reply}

\begin{objection}[Normative power]
On the normative power view, a promise binds because the promisor communicates his intention to bind himself, and the promisee's interest in having authority over the promisor's conduct explains why there should be such a power \citep{Owens2012}. Reliance plays no role, and the persistence of the obligation after forgetting is explained without it.
\end{objection}

\begin{reply}
The normative power view explains persistence well, and our account agrees with it there. It explains insincere promises less well, since the promisor who does not intend to bind himself does not exercise the power, and yet he is bound. Our account explains both: the power, if we call it that, is exercised by inviting reliance and presenting oneself as answerable for it, whatever one's private intention; and the obligation so created is grounded in a past act, whose effects outlast the promisee's attention.
\end{reply}

\section{What the result does not establish}

The argument does not determine how stringent promissory obligations are, or when they are defeated by weightier reasons; a promise to do something seriously wrong creates no obligation to do it, and promises can be outweighed, as ETH-15 and ETH-27 discuss. Nor does it settle whether promises to the dead bind, which raises the further question whether the dead can have entitlements. It shows only that, among living promisees, actual reliance is neither necessary for the creation of the obligation nor necessary for its persistence.

\section{Conclusion}

Promissory obligation arises from the creation of reasonable reliance: the promisor invites the promisee to rely on him and presents himself as answerable for her doing so, and this entitles her to rely and obligates him to perform. Because the obligation is grounded in that act and directed to the promisee's entitlement, it persists until it is discharged or released. Forgetting is neither, and so the promise binds the promisor even when the promisee no longer remembers it.

\begin{thebibliography}{99}\small
\bibitem[Hume(1739)]{Hume1739} Hume, D. (1739--1740). \emph{A Treatise of Human Nature}. London: John Noon. Cited by book, part and section.
\bibitem[Kolodny and Wallace(2003)]{KolodnyWallace2003} Kolodny, N. and Wallace, R.~J. (2003). Promises and practices revisited. \emph{Philosophy \& Public Affairs}, 31(2), 119--154.
\bibitem[Owens(2012)]{Owens2012} Owens, D. (2012). \emph{Shaping the Normative Landscape}. Oxford University Press.
\bibitem[Rawls(1955)]{Rawls1955} Rawls, J. (1955). Two concepts of rules. \emph{The Philosophical Review}, 64(1), 3--32.
\bibitem[Raz(1977)]{Raz1977} Raz, J. (1977). Promises and obligations. In P.~M.~S. Hacker and J. Raz (eds), \emph{Law, Morality, and Society: Essays in Honour of H.~L.~A. Hart}. Oxford: Clarendon Press.
\bibitem[Scanlon(1990)]{Scanlon1990} Scanlon, T.~M. (1990). Promises and practices. \emph{Philosophy \& Public Affairs}, 19(3), 199--226.
\bibitem[Scanlon(1998)]{Scanlon1998} Scanlon, T.~M. (1998). \emph{What We Owe to Each Other}. Cambridge, MA: Harvard University Press.
\bibitem[Searle(1969)]{Searle1969} Searle, J.~R. (1969). \emph{Speech Acts: An Essay in the Philosophy of Language}. Cambridge University Press.
\bibitem[Shiffrin(2008)]{Shiffrin2008} Shiffrin, S.~V. (2008). Promising, intimate relationships, and conventionalism. \emph{The Philosophical Review}, 117(4), 481--524.
\end{thebibliography}

\end{document}
